% Minimal class-compatibility preamble for a reader excerpt of
%   Frederic Deglise, Jean Fasel, Fangzhou Jin, Adeel A. Khan,
%   "On the rational motivic homotopy category",
%   Journal de l'Ecole polytechnique -- Mathematiques 8 (2021), DOI 10.5802/jep.153.
%
% The block below is the author's own macro, package and theorem declaration block of the
% pinned source file, copied mechanically (source lines 6-16 and 18-21 and 37-135). Only the
% publisher's OT2/cyrillic font block (source lines 22-35) is left out, because it declares
% 8-bit cyrillic fonts that the retained text never uses; a Unicode text font is declared
% instead so that the accented letters of the retained author text are rendered. No
% mathematical macro, statement, symbol or assumption is introduced, changed or removed.
\usepackage{fontspec}
\setmainfont{lmroman10-regular.otf}[
  ItalicFont=lmroman10-italic.otf,
  BoldFont=lmroman10-bold.otf,
  BoldItalicFont=lmroman10-bolditalic.otf]
\usepackage{xspace}

\def\To#1{\mathchoice{\xrightarrow{\textstyle\kern4pt#1\kern3pt}}{\stackrel{#1}{\longrightarrow}}{}{}}
\def\mto{\mathchoice{\longmapsto}{\mapsto}{\mapsto}{\mapsto}}
\newcommand{\Psfrac}[2]{(\sfrac{#1}{#2})}
\newcommand{\psfrac}[2]{\sfrac{(#1)}{#2}}
\newcommand{\spfrac}[2]{\sfrac{#1}{(#2)}}
\newcommand{\pspfrac}[2]{\sfrac{(#1)}{(#2)}}
\let\epsilon\varepsilon

\newcommand{\tstructure}{$t$-structure\xspace}
\newcommand{\tstructures}{$t$-structures\xspace}
\newcommand{\BM}{\mathrm{BM}}
\usepackage[all,2cell,emtex,cmtip]{xy}\let\labelstyle\textstyle
\usepackage{bbm,nicefrac}
\usepackage[mathscr]{euscript}
\usepackage{mathrsfs}
\newcommand{\modl}[1]{#1\textnormal{-}\mathrm{mod}}
\newcommand{\base}{\mathscr S}

\newcommand{\T}{\mathscr T}
\newcommand{\C}{\mathscr C}

\DeclareMathOperator{\uK}{\underline K}
\DeclareMathOperator{\uPic}{\underline{Pic}}
\DeclareMathOperator{\uZZ}{\underline{\ZZ}}
\DeclareMathOperator{\rk}{rk}
\DeclareMathOperator{\GW}{GW}

\renewcommand{\AA}{\mathbb A}
\newcommand{\E}{\mathbb E}
\newcommand{\F}{\mathbb F}
\newcommand{\un}{\mathbbm 1}
\newcommand{\MSL}{\mathbf{MSL}}
\newcommand{\KQ}{\mathbf{KQ}}
\newcommand{\KW}{\mathbf{KW}}
\newcommand{\KMW}{\mathbf K^\mathrm{MW}}
\newcommand{\uW}{\underline{\mathrm W}}
\newcommand{\HuW}{\mathbf H \underline{\mathrm W}}
\newcommand{\HM}{\mathbf H_{\mathrm{M}}}
\newcommand{\HMW}{\mathbf H_{\mathrm{MW}}}
\newcommand{\HW}{\mathbf H_{\mathrm{W}}}
\newcommand{\HH}{\mathrm H}
\newcommand{\HA}{\mathrm H_{\AA^1}}
\newcommand{\BA}{\mathrm H^{\AA^1}}

\newcommand{\GL}{\mathrm{GL}}
\newcommand{\SL}{\mathrm{SL}}

\newcommand{\thom}{\mathfrak t}
\newcommand{\vb}[1]{\langle#1\rangle}

\DeclareMathOperator{\codim}{codim}
\newcommand{\eff}{\mathrm{eff}}

\newcommand{\sm}{\mathrm{Sm}}
\DeclareMathOperator{\DM}{DM}
\DeclareMathOperator{\DMB}{DM_\mathcyr B}
\DeclareMathOperator{\SH}{SH}
\DeclareMathOperator{\DA}{D_{\AA^1}}
\DeclareMathOperator{\DAeff}{D_{\AA^1}^{\eff}}
\DeclareMathOperator{\DAminus}{D_{\AA^1,-}}
\DeclareMathOperator{\Id}{Id}
\DeclareMathOperator{\Hom}{Hom}
\DeclareMathOperator{\Maps}{Maps}
\DeclareMathOperator{\uHom}{\underline{Hom}}
\DeclareMathOperator{\Pic}{Pic}
\DeclareMathOperator{\CH}{CH}
\DeclareMathOperator{\wCH}{\widetilde{CH}}
\DeclareMathOperator{\spec}{Spec}
\DeclareMathOperator{\Ind}{Ind}
\DeclareMathOperator{\ilim}{\varinjlim}
\DeclareMathOperator{\plim}{\varprojlim}
\newcommand{\twist}[1]{\langle #1 \rangle}
\DeclareMathOperator{\Th}{Th}
\DeclareMathOperator{\Tw}{Tw}
\DeclareMathOperator{\tdet}{\widetilde \det}
\DeclareMathOperator{\tomega}{\widetilde \omega}
\DeclareMathOperator{\K}{K}
\DeclareMathOperator{\Der}{D}
\DeclareMathOperator{\Spt}{Spt}

\newcommand{\ZZ} {\mathbb Z}
\newcommand{\ZZot} {\mathbb Z[\nicefrac 1 2]}
\newcommand{\ot} {{\nicefrac 1 2}}
\newcommand{\QQ} {\mathbb Q}
\newcommand{\OO}{\mathcal O}
\newcommand{\PP} {\mathbb P}
\newcommand{\GG} {\mathbb{G}_m}
\newcommand{\GGx}[1] {\mathbb{G}_{m,#1}}
\newcommand{\nis}{\mathrm{Nis}}
\newcommand{\zar}{\mathrm{Zar}}

\newtheorem{thm}{Theorem}[section]
\newtheorem{prop}[thm]{Proposition}
\newtheorem{lm}[thm]{Lemma}
\newtheorem{cor}[thm]{Corollary}
\newtheorem{thmi}{Theorem}
\newtheorem{cori}[thmi]{Corollary}
\newtheorem*{lmi}{Key Lemma}

\renewcommand{\thethmi}{\Alph{thmi}}

\theoremstyle{remark}
\newtheorem{rem}[thm]{Remark}
\newtheorem{notation}[thm]{Notation}
\newtheorem{ex}[thm]{Example}
\newtheorem{conv}[thm]{Convention}
\newtheorem{setup}[thm]{Setup}
\newtheorem{warn}[thm]{Warning}

\theoremstyle{definition}
\newtheorem{df}[thm]{Definition}

\let\thesubsubsection\thethm
\newcommand{\numer}[1]{\refstepcounter{thm}\subsubsection{#1}}
\newtheorem{num}[thm]{}
\numberwithin{equation}{thm}

\renewcommand{\theequation}{\thethm.\alph{equation}}

% Publisher-class abbreviations (cedram) used by the retained author text.
\providecommand{\loccit}{\emph{loc. cit.}}
\providecommand{\bibliofont}{\small}
\providecommand{\backmatter}{}
\providecommand{\skpt}{\smallskip}
\providecommand{\smaller}{\small}
\newcommand{\ie}{i.e.,\xspace}
\newcommand{\eg}{e.g.,\xspace}
\newcommand{\etc}{etc.\xspace}
\newcommand{\cf}{cf.\xspace}
\newcommand{\resp}{resp.\xspace}
